% !TEX program = pdflatex
% Generic preprint skeleton on the standard `article` class -- BibLaTeX + Biber.
%
% IMPORTANT: This is a NEUTRAL preprint layout, NOT a substitute for any venue
% template. When you target IEEE, ACM, Springer, Elsevier, etc., download and
% use that venue's OFFICIAL class (IEEEtran, acmart, sn-jnl, elsarticle, ...)
% and move your content into it. Do not ship this layout to a venue that
% mandates its own.
%
% Build:  latexmk main.tex
% Clean:  latexmk -c
%
% For a two-column preprint, add `twocolumn` to the class options:
%   \documentclass[11pt,a4paper,twocolumn]{article}
\documentclass[11pt,a4paper]{article}

% --- Encoding & fonts (pdfLaTeX) --------------------------------------------
\usepackage[T1]{fontenc}
\usepackage[utf8]{inputenc}

% --- Core packages ----------------------------------------------------------
\usepackage[margin=1in]{geometry}
\usepackage{microtype}
\usepackage{amsmath}
\usepackage{amssymb}
\usepackage{graphicx}
\usepackage{booktabs}

% --- Bibliography: BibLaTeX with the Biber backend --------------------------
\usepackage[backend=biber,style=numeric,sorting=nyt]{biblatex}
\addbibresource{refs.bib}

% --- Hyperlinks & smart cross-references ------------------------------------
\usepackage{hyperref}
\usepackage{cleveref}

\graphicspath{{figures/}}

% A lightweight "keywords" command (venue classes usually provide their own).
\providecommand{\keywords}[1]{\par\noindent\textbf{Keywords:} #1}

\title{A Generic Preprint Skeleton}
\author{%
  Author One\\
  \small Department of Placeholder Studies, Placeholder University\\
  \small \texttt{author.one@example.org}
  \and
  Author Two\\
  \small Placeholder Research Lab\\
  \small \texttt{author.two@example.org}
}
\date{\today}

\begin{document}
\maketitle

\begin{abstract}
Placeholder abstract. State the problem, the approach, and the main result in a
few sentences suitable for a preprint server.
\end{abstract}

\keywords{placeholder, template, preprint, LaTeX}

\section{Introduction}
\label{sec:intro}
Placeholder introduction. Cite prior work such as~\cite{placeholder-paper} and
cross-reference \Cref{eq:model}, \Cref{tab:results}, and \Cref{fig:overview}.

\section{Method}
\label{sec:method}
A display equation:
\begin{equation}
  \label{eq:model}
  \hat{y} = f_\theta(x).
\end{equation}

\begin{figure}[htbp]
  \centering
  \includegraphics[width=0.7\linewidth]{example.pdf}
  \caption{Placeholder overview figure. Put \texttt{example.pdf} in
    \texttt{figures/}.}
  \label{fig:overview}
\end{figure}

\section{Results}
\label{sec:results}

\begin{table}[htbp]
  \centering
  \caption{Placeholder results using \texttt{booktabs} rules.}
  \label{tab:results}
  \begin{tabular}{lrr}
    \toprule
    Method   & Metric A & Metric B \\
    \midrule
    Baseline & 0.81     & 12.4     \\
    Ours     & 0.89     &  9.7     \\
    \bottomrule
  \end{tabular}
\end{table}

\section{Conclusion}
\label{sec:conclusion}
Placeholder conclusion.

\printbibliography

\end{document}
